\section{Introduction}
Understanding a painting is not a single act of perception. It is a layered process that draws simultaneously on visual attention, affective response, formal knowledge, and historical memory. A viewer encountering Vermeer's Girl with a Pearl Earring may first register the luminous skin and direct gaze — a narrative observation. A trained eye may then notice the sfumato technique and restricted palette — a formal judgment. The same viewer may feel unease or intimacy — an emotional response. And an art historian will situate the work within the Dutch Golden Age domestic portrait tradition — a contextual interpretation. These four acts of understanding are not sequential stages but concurrent, interacting perspectives on a single visual object.

Current vision-language models do not work this way. Despite rapid advances in multimodal understanding, existing systems are predominantly trained on object-centric, single-perspective captions that flatten this interpretive complexity into a single descriptive register. Benchmarks and datasets in computational art analysis — SemArt, ExpArt, ArtEmis, OmniArt — have each captured important dimensions of art understanding, but none provides a unified, large-scale resource that systematically represents all four interpretive dimensions simultaneously for the same artworks. As a result, models trained on these resources learn to describe paintings, but not to interpret them in the full sense that word carries in art history and criticism.

We introduce PolyArt, a large-scale multi-perspective dataset of over 70,000 paintings annotated across four interpretive dimensions: Narrative and Scene Interpretation, Formal Visual Analysis, Emotional Response, and Historical and Contextual Analysis. Each painting in PolyArt is paired with four independently generated captions, one per perspective, produced through a hybrid annotation pipeline combining state-of-the-art multimodal models including Qwen2.5-VL, GalleryGPT, ArtRAG, and ArtEmis. The perspectives are generated independently to preserve interpretive diversity and merged into a unified polysemic caption through an explanation harmonization step. PolyArt is the first dataset to provide this breadth of interpretive coverage at this scale, with open licensing and reproducible generation protocols.

Beyond its value as a descriptive resource, PolyArt is designed as an evaluation environment for a question that has received little systematic attention: does the reasoning mode an agent employs when interpreting art affect the quality and character of the understanding it produces? Generating a formal analysis of a painting requires different cognitive operations than generating an emotional response or a historical contextualisation — operations that correspond, in the agentic AI literature, to distinct reasoning strategies including analogical reasoning, causal inference, semiotic interpretation, and phenomenological description. PolyArt makes these distinctions tractable: because each perspective is generated and evaluated independently, it becomes possible to measure whether agentic systems that are prompted to reason in perspective-appropriate ways outperform those given only generic generation instructions.

We demonstrate this through two experimental contributions that extend the dataset into an active benchmark. First, we introduce a perspective sufficiency analysis using image regeneration: given only the textual perspectives of a painting — without the original image — we test whether generative models can reconstruct the visual content, and we measure which perspective subsets are sufficient to recover which dimensions of the original (style, composition, iconographic content). This provides an empirical grounding for the claim that PolyArt's perspectives encode distinct and complementary information rather than paraphrasing each other. Second, we conduct a reasoning mode evaluation in which agentic systems are prompted with different reasoning strategies — analogical, causal, formal, and counterfactual — across perspective-specific generation tasks, revealing systematic variation in output quality and interpretive accuracy depending on the alignment between reasoning mode and perspective type.

Together, these contributions position PolyArt not merely as a larger or richer captioning dataset, but as a resource for studying how multimodal agents reason about visual culture — a question with implications beyond the art domain for the broader development of evaluable, interpretable agentic AI systems.


\noindent Our contributions are:
\begin{itemize}
    \item \textbf{Multi-Perspective Dataset.} We introduce PolyArt, the first large-scale multimodal dataset providing four independently generated interpretive perspectives for approximately 70,000 paintings, covering narrative, formal, emotional, and historical dimensions simultaneously for the same artworks. Each painting additionally receives a harmonized unified caption synthesizing all four perspectives.

    \item \textbf{Specialized Generation Pipeline.} We design a heterogeneous pipeline that assigns each perspective to the vision-language model best suited for the corresponding interpretive task --- Qwen2.5-VL-72B for narrative and historical analysis, GalleryGPT for formal analysis, and Qwen3-VL-8B grounded in ArtEmis human utterances for emotional response --- maximizing per-perspective quality rather than treating annotation as a homogeneous generation task.

    \item \textbf{Perspective Complementarity Benchmark.} We operationalize the complementarity of perspectives through an image reconstruction experiment: given text from one or more perspective subsets, a generative model attempts to reconstruct the original painting, and fidelity is measured across style, composition, iconographic content, and emotional tone. This provides a falsifiable empirical test of whether the perspectives encode distinct visual information.

    \item \textbf{Reasoning Mode Alignment Benchmark.} We introduce a reasoning mode evaluation in which agentic systems are prompted with six perspective-appropriate reasoning strategies and assessed on perspective-targeted generation tasks, revealing whether alignment between reasoning mode and interpretive dimension produces systematically better explanations than generic prompting.
\end{itemize}